\minitopic{拥有理由与基于理由}\term{命题确证}{propositional justification}指主体拥有足以支持$p$的根据；\term{信念确证}{doxastic justification}还要求主体因为这些根据而相信$p$。开篇中甲乙或许都“有”同一病例资料，但乙的信念没有适当基于它。这个区别防止我们把偶然拥有未使用证据当作认知成就。 证据也可能被\term{击败}{defeat}。看到钟指向九点通常支持“现在九点”，但得知钟已停摆便产生反驳型击败；得知房间布满外观相同的坏钟，则削弱证据与结论的联系。理性主体不仅累积支持，还须对击败者敏感。

\minitopic{回溯与结构}若每个得到确证的信念都需要另一个信念支持，就出现无穷回溯。\term{基础主义}{foundationalism}以某些非推论信念终止回溯，典型候选是当前经验或内省。基础信念不必不可错；温和基础主义只要求它们拥有无需其他信念传递的初始确证。 \term{融贯主义}{coherentism}认为确证来自信念系统的相互支持、一致性和解释整合\parencite{bonjour1985}。它避免单向地基，却面对“孤立系统”与“多个同样融贯系统”问题：一部内部一致的虚构故事为何没有世界知识？融贯必须加入经验输入或可靠联系。

\minitopic{内在主义与外在主义}内在主义要求决定确证的因素以某种方式处于主体的反思可及范围。它保护第一人称责任：主体应能说明自己为何相信。外在主义则允许主体不知道使其信念可靠的全部事实；儿童和动物可通过正常知觉知道，而无需掌握关于视觉系统的理论。 \term{证据主义}{evidentialism}主张信念态度应与主体证据相称\parencite{conee1985}。\term{可靠主义}{reliabilism}则把信念确证联系到产生信念的可靠过程\parencite{goldman1979}。前者面对证据究竟由什么构成的问题，后者面对“过程类型如何划分”以及新邪恶恶魔问题：一个受系统欺骗但内部经验与我们相同的主体，似乎仍然理性，即使其过程在环境中不可靠。

\begin{argumentbox}
可靠主义的动机：可靠知觉与不可靠猜测的差异不必被主体完全表示；知识面向真理，因此产生真信念比例高的过程具有认识意义。反驳不能只说主体“不知道自己可靠”，因为这正是假定内在主义；它需要说明反思可及为何是确证不可缺少的条件。
\end{argumentbox}

适当功能理论进一步要求认知能力在适宜环境中依照其设计方案运作并以真理为目标\parencite{plantinga1993}。它把可靠性置入更广的功能框架，但“设计”“正常环境”和目标的解释均有争议。

\begin{comparison}
朱熹“格物穷理”与王阳明“致良知”的争论有外向与内向的认识路径差异，却不是可靠主义与内在主义的历史版本\parencite{tiwald2014}。前者嵌在理、心与工夫论中，目标包括道德成就。恰当比较应问：根据如何成为主体的根据？反思、实践和环境分别承担什么角色？
\end{comparison}

\minitopic{确证的时间维度}信念可能在形成时得到确证，后来因新信息失去确证；也可能起初只是猜测，后来获得支持。认识评价必须注明时间。主体收到反证却从未注意，是已经失去确证，还是只在注意后才失去？心理主义答案要求击败者进入主体心理状态；规范事实主义答案可能认为主体可得却被疏忽的信息已经构成问题。 还有\term{高阶证据}{higher-order evidence}：它不直接支持或反对$p$，而是评价主体处理一阶证据的能力。例如得知自己服用了影响计算的药物，并不直接证明答案错误，却削弱对计算结果的信任。成熟确证理论必须说明一阶证据与关于自身可靠性的证据如何互动。

\minitopic{理由的本体是什么}有人把理由理解为主体的心理状态，如经验或信念；有人认为理由是事实；还有人把证据等同于主体所知道的命题\parencite{williamson2000}。事实观解释理由为何真正支持结论，却面对主体不知道的事实如何成为“他的”理由。心理状态观保证拥有关系，却可能允许完全相同的幻觉经验像真知觉一样提供理由。知识观兼顾事实性与拥有关系，但会使知识参与解释确证，改变传统分析顺序。 选择理由本体会影响案例判断。正常看见红灯与完美幻觉在内部相同：若理由是经验，两者拥有相同初步理由；若理由是“灯是红的”这个已知事实，只有前者拥有该理由。争论不能只停在“他们看起来同样合理”，还要说明所评价的是内部无责、客观支持还是知识资格。

\minitopic{案例工作坊：研究数据}研究者正确运行统计程序，得到显著结果；同事随后指出数据文件标签可能调换。此时原结果受到击败。若研究者检查日志确认标签无误，击败者被解除；若他只因不喜欢质疑而忽略邮件，信念即使碰巧为真，也缺少适当维持。 这个案例区分信念的产生与维持。可靠过程不能只评价最初计算，还要包括响应反证的更新机制。认识责任也不是要求对每个想象可能无限复核，而是对具体、相关且来自可信渠道的挑战作出相称响应。

\minitopic{理论组合而非标签选择}基础主义者可以是内在主义者，也可以把基础信念建立在外在可靠关系上；融贯主义者也可要求系统由可靠输入约束。证据主义与可靠主义有时竞争，有时分别评价内部合宜与知识导向。学习这些理论的目标不是选择一个阵营名称，而是明确每个理论解释哪种认识评价。 一种多元框架可同时承认：经验给予初步理由，信念需与整体证据融贯，形成过程需在环境中可靠，主体还应响应可得击败者。代价是理论不再用单一原则解释一切；优点是更贴近日常评价的多维结构。
